\documentclass[11pt]{article}
\ifdefined\pdfmapfile\pdfmapfile{+lm.map}\pdfmapfile{+cm.map}\pdfmapfile{+symbols.map}\fi
\usepackage[T1]{fontenc}
\usepackage{lmodern}
\usepackage[margin=1in]{geometry}
\usepackage{amsmath,amssymb,amsthm,hyperref}
\hypersetup{hidelinks}
\newtheorem{theorem}{Theorem}
\newtheorem{corollary}[theorem]{Corollary}
\newtheorem{lemma}[theorem]{Lemma}
\newcommand{\qbin}[2]{\genfrac{[}{]}{0pt}{}{#1}{#2}_q}
\newcommand{\fmaj}{\operatorname{fmaj}}
\newcommand{\Des}{\operatorname{Des}}
\title{Gaussian Products and Hyperoctahedral Descent Classes:\\A Parabolic-Quotient Proof of Unimodality}
\author{Research proof note}
\date{November 2026}
\begin{document}
\maketitle
\begin{abstract}
We give a complete proof of the statements printed as Conjectures 5.1--5.3
in the 2005 preprint of Adin, Brenti, and Roichman. The central polynomial
$\qbin nk\prod_{j=k+1}^n(1+q^j)$ is the rank polynomial of the Weyl-group
quotient of type $B_n/(B_k\times A_{n-k-1})$. Stanley's classical hard
Lefschetz theorem for parabolic quotients therefore proves symmetry and
unimodality in every rank. We make the reduction of the endpoint-refined
statements explicit, including empty classes and monomial shifts, and give
an interpretation of the successive coefficient differences as primitive
cohomology dimensions. The argument is an application of established theory;
no claim that this observation is historically new or that the printed
conjectures have remained unresolved since 2005 is made.
\end{abstract}

\section{Statement and conventions}
Set $[m]_q=1+q+\cdots+q^{m-1}$ for $m\ge1$, $[0]_q!=1$, and
$[m]_q!=\prod_{j=1}^m[j]_q$. For nonnegative $d_1,\ldots,d_t$ with
sum $N$, put
\[
\qbin N{d_1,\ldots,d_t}=\frac{[N]_q!}{\prod_{s=1}^t[d_s]_q!}.
\]
A nonzero polynomial is called symmetric on its support if it has the form
$q^aP(q)$ where $P(0)\ne0$ and $P(q)=q^{\deg P}P(q^{-1})$.
Unimodality allows weak inequalities. The zero polynomial is regarded as
unimodal. This support convention matters for fixed endpoints.

The statement printed as Conjecture~5.3 of~\cite{ABR} is the following.
\begin{theorem}\label{thm:main}
For every pair of integers $0\le k\le n$, the polynomial
\[
F_{n,k}(q)=\qbin nk\prod_{j=k+1}^n(1+q^j)
\]
has nonnegative integer coefficients and is symmetric and unimodal.
Its degree and value at one are
\[
D=\frac{(n-k)(n+3k+1)}2,\qquad F_{n,k}(1)=2^{n-k}\binom nk.
\]
The case $k=n$ is included by the convention $F_{n,n}=1$.
\end{theorem}

\section{The parabolic quotient}
We use the following established theorem, in the precise form proved by
Stanley~\cite[Theorem~3.1]{Stanley}.
\begin{lemma}[Stanley]
If $W$ is a finite Weyl group and $W_J$ a standard parabolic subgroup,
then the length-generating polynomial of the minimal coset representatives
$W^J$ is symmetric and unimodal.
\end{lemma}
For completeness, the geometric implication in this result is as follows.
The quotient $W^J$ indexes the Schubert cells of the smooth projective
homogeneous variety $X=G/P$. The cells have complex dimensions $\ell(w)$,
so their length polynomial is
\[
\sum_{w\in W^J}q^{\ell(w)}
 =\sum_{i=0}^{\dim_{\mathbb C}X}\dim H^{2i}(X;\mathbb Q)q^i.
\]
If $L$ is multiplication by an ample class, hard Lefschetz gives
\[
L^{D-2i}:H^{2i}(X;\mathbb Q)\longrightarrow
H^{2D-2i}(X;\mathbb Q)
\quad\text{an isomorphism }(0\le2i\le D).
\]
It follows that $L:H^{2i}\to H^{2i+2}$ is injective for $i<D/2$.
Poincare duality supplies symmetry, proving the lemma's two conclusions.

\begin{proof}[Proof of Theorem~\ref{thm:main}]
The finite Weyl-group length polynomials are
\[
P_{B_n}(q)=\prod_{j=1}^n[2j]_q,
\qquad P_{S_r}(q)=[r]_q!.
\]
Choose the standard parabolic subgroup with components $B_k$ and
$A_{n-k-1}$, with the usual convention that a component of rank zero is
trivial. Its Weyl group is $B_k\times S_{n-k}$. Every $w\in B_n$ has a
unique parabolic factorization $w=uv$, with $u\in W^J$, $v\in W_J$, and
$\ell(w)=\ell(u)+\ell(v)$. Consequently,
\begin{align*}
\sum_{u\in W^J}q^{\ell(u)}
 &=\frac{\prod_{j=1}^n[2j]_q}
 {\bigl(\prod_{j=1}^k[2j]_q\bigr)[n-k]_q!}\\
 &=\frac{\prod_{j=k+1}^n[j]_q(1+q^j)}{[n-k]_q!}\\
 &=\qbin nk\prod_{j=k+1}^n(1+q^j)=F_{n,k}(q).
\end{align*}
Stanley's theorem now applies directly. Counting degrees gives
\[
\sum_{j=k+1}^n(2j-1)-\frac{(n-k)(n-k-1)}2
 =\frac{(n-k)(n+3k+1)}2.
\]
Evaluation at $q=1$ gives the asserted value. If $k=n$, the conclusion
is immediate and no maximal parabolic subgroup is needed.
\end{proof}

One concrete geometric model is the symplectic isotropic Grassmannian
$\operatorname{IG}(n-k,2n)$. Its parabolic Levi factor is
$\operatorname{GL}_{n-k}\times\operatorname{Sp}_{2k}$, with Weyl group
$S_{n-k}\times B_k$. Its dimension is precisely $D$ above. Thus no singular
variety or conjectural extension of hard Lefschetz is being used.

\section{The two descent-class statements}
Write $B_n$ as the signed permutations of $\{1,\ldots,n\}$, with
$\sigma(-j)=-\sigma(j)$. Use the Coxeter length $\ell_B$ and
\[
\Des_B(\sigma)=\{j\in\{0,\ldots,n-1\}:\sigma(j)>\sigma(j+1)\},
\qquad\sigma(0)=0.
\]
For $M\subseteq\{0,\ldots,n-1\}$ and $i\in\{\pm1,\ldots,\pm n\}$,
define
\[
G_M(q)=\sum_{\Des_B(\sigma^{-1})\subseteq M}q^{\ell_B(\sigma)},
\qquad
G_{M,i}(q)=\sum_{\substack{\Des_B(\sigma^{-1})\subseteq M\\\sigma(n)=i}}
q^{\ell_B(\sigma)}.
\]
These are exactly the length versions of Conjectures~5.2 and~5.1
in~\cite{ABR}. Their equality with the corresponding flag-major-index
polynomials is already proved in Theorems~3.2 and~3.3 of that paper.

\begin{corollary}\label{cor:classes}
For all $n$, $M$, and $i$ as above, $G_M(q)$ is symmetric and unimodal,
and $G_{M,i}(q)$ is either zero or symmetric on its support and unimodal.
The same statements hold with $\ell_B$ replaced by $\fmaj$.
\end{corollary}
\begin{proof}
If $M=\varnothing$, the only permitted permutation is the identity.
Thus $G_M=1$, $G_{M,n}=1$, and $G_{M,i}=0$ for $i\ne n$.
Assume $M=\{m_1<\cdots<m_t\}$ is nonempty. Put
\[
k=m_1,\qquad m_{t+1}=n,\qquad d_s=m_{s+1}-m_s\quad(1\le s\le t).
\]
Then $d_s\ge1$ and $\sum_s d_s=n-k$. The product formula of
\cite[Theorem~3.3]{ABR} gives
\begin{equation}\label{eq:GM}
G_M(q)=F_{n,k}(q)\qbin{n-k}{d_1,\ldots,d_t}.
\end{equation}
A Gaussian multinomial is a product of Gaussian binomials and is symmetric
and unimodal; alternatively it is the length polynomial of a type-$A$
parabolic quotient. Products of nonnegative symmetric unimodal
polynomials are symmetric and unimodal. Thus~\eqref{eq:GM} proves the
unrefined assertion.

For clarity, the endpoint formulas from~\cite[Theorems~4.2--4.3]{ABR}
can be rewritten without any rational-function cancellation:
\begin{align}
G_{M,k}(q)
 &=q^{n-k}F_{n-1,k-1}(q)\qbin{n-k}{d_1,\ldots,d_t}
 &&(k>0),\label{eq:ep1}\\
G_{M,m_{s+1}}(q)
 &=q^{n-m_{s+1}}F_{n-1,k}(q)
 \qbin{n-k-1}{d_1,\ldots,d_s-1,\ldots,d_t},\label{eq:ep2}\\
G_{M,-(m_s+1)}(q)
 &=q^{n+m_s}F_{n-1,k}(q)
 \qbin{n-k-1}{d_1,\ldots,d_s-1,\ldots,d_t}
 &&(1\le s\le t).\label{eq:ep3}
\end{align}
All other endpoint polynomials vanish. Zero entries in a multinomial
are allowed, and $F_{a,a}=1$ covers the endpoint case $k=n-1$.
These identities follow directly on substituting the three values of
$\alpha_i(M)$ in~\cite[Theorem~4.2]{ABR}, using
$[2n]_q=[n]_q(1+q^n)$ and cancelling the corresponding factorial factor.
Each right-hand side is a monomial times a product of the already proved
symmetric unimodal polynomials. This proves the assertion. Finally,
the established equidistribution identities transfer it to $\fmaj$.
\end{proof}

For reference, the product fact used above has a short elementary proof.
Every nonnegative symmetric unimodal polynomial of degree $D$ has a
unique expansion
\[
\sum_{j=0}^{\lfloor D/2\rfloor}b_jq^j[D-2j+1]_q,\qquad b_j\ge0.
\]
The coefficients of $[a]_q[b]_q$ increase, may plateau, and decrease
symmetrically. Each product of two summands has the same symmetry center
as the full product. Their nonnegative sum is therefore unimodal.

\section{Primitive decomposition and scope}
\begin{corollary}
Write $F_{n,k}(q)=\sum_{i=0}^D a_iq^i$, and set $a_{-1}=0$. Then
\[
b_i=a_i-a_{i-1}\in\mathbb Z_{\ge0}
\quad(0\le i\le\lfloor D/2\rfloor),
\qquad
F_{n,k}(q)=\sum_{i=0}^{\lfloor D/2\rfloor}b_iq^i[D-2i+1]_q.
\]
Moreover, $b_i$ is the dimension of the primitive cohomology in degree
$2i$ of $\operatorname{IG}(n-k,2n)$, with respect to an ample class.
\end{corollary}
\begin{proof}
Hard Lefschetz and its primitive decomposition give independent strings
$u,Lu,\ldots,L^{D-2i}u$ beginning in degree $2i$ for primitive classes $u$.
Each such string contributes $q^i[D-2i+1]_q$. Comparing consecutive
coefficients before the middle gives $b_i=a_i-a_{i-1}$.
\end{proof}
For example,
\[
F_{3,1}(q)=1+q+2q^2+2q^3+2q^4+2q^5+q^6+q^7
 =[8]_q+q^2[4]_q.
\]

The conclusions concern $\Des_B(\sigma^{-1})\subseteq M$. Replacing
this condition by equality is a different statement and is not established
by this argument. Likewise, rank unimodality does not by itself produce
a combinatorial symmetric-chain decomposition.

The exact source of the conjecture labels is arXiv:math/0508362v4,
Section~5.1 (2005); the paper appeared in 2006. The main theorem above
is an explicit application of Stanley's 1980 theorem. A limited literature
search did not locate a separately named later resolution of the three
printed conjectures, but this does not establish priority, continuing
open status, or a new theorem beyond classical parabolic-quotient theory.

\begin{thebibliography}{9}
\bibitem{ABR}
R.~M.~Adin, F.~Brenti, and Y.~Roichman,
\emph{Equi-distribution over descent classes of the hyperoctahedral group},
J. Combin. Theory Ser. A \textbf{113} (2006), 917--933.
\href{https://doi.org/10.1016/j.jcta.2005.08.008}{doi:10.1016/j.jcta.2005.08.008}.
The conjecture numbering refers to the 2005 preprint, version~4:
\url{https://arxiv.org/abs/math/0508362}.
\bibitem{Stanley}
R.~P.~Stanley,
\emph{Weyl groups, the hard Lefschetz theorem, and the Sperner property},
SIAM J. Algebraic Discrete Methods \textbf{1} (1980), 168--184.
Theorem~3.1.
\href{https://doi.org/10.1137/0601021}{doi:10.1137/0601021}.
Author-hosted full text:
\url{https://math.mit.edu/~rstan/pubs/pubfiles/42.pdf}.
\end{thebibliography}
\end{document}
